\section{Introduction}
\label{sec:introduction}
Voice assistants engage with users through direct, live speech-to-speech models such as Gemini Live \cite{google2026geminilive}. With the appearance of more voice based applications, the ability to engage with these tools in acoustically challenging environments is becoming increasingly important. A particular problem is the presence of interfering speech, which is misinterpreted by the live speech models as the intended target user speech. Interfering speech thus will confuse the downstream model as to the intended user input, which will lead to a poor user experience. Secondly, live speech models are exposed to noisy environments that include reverberation, background audio and other acoustic challenges. Although live models are fairly robust to background noise on its own (\autoref{sec:llms}), processing audio to remove interfering speech can introduce artefacts that the live model is not robust to. The work in this report aims to address these issues through the development of a target speaker extractor that aims to preprocess the audio input that is given to the live model. The extractor is constructed under the following constraints: it must be causal, process audio in real time and operate on a single channel of audio. However, an extractor can itself introduce artefacts that the live model hears as words, so it is not enough to measure whether the interfering speaker was removed; the evaluation must show which of the two errors the extractor causes.

Under these constraints, the contribution of this work lies not in a new architecture but in how the extractor is evaluated. The novel contributions are:
\begin{enumerate}
    \item \textbf{A live-model content fidelity evaluation suite.} The live model transcribes the extractor's output, and its errors are split into words leaked from the interfering speaker and words that neither speaker said (\autoref{sec:metrics}).
    \item \textbf{Evidence that the metrics are needed.} A single word error rate hides changes that the split reveals, and an offline transcriber misjudges changes that the live model does not notice (\autoref{sec:results}).
\end{enumerate}
To produce and test these, the following was built:
  \begin{enumerate}
    \setcounter{enumi}{2}
    \item \textbf{A constructed dataset and evaluation protocol.} \num{9955} training trials from \num{1172} LibriSpeech speakers, each rendered by constructing a simulated room with a random reverberation time and microphone and speaker positions, and then adding real background noise. Each speaker has exact transcripts. The dataset includes trials where the target is absent so that the extractor learns to stay silent (\autoref{sec:data}). 
    \item \textbf{A streaming target speaker extractor.} The REAL-TSE baseline architecture \cite{wang2026realtse}, implemented for this work and made fully causal, with the enrollment cue split into parts and a speaker embedding added (\autoref{sec:architecture}, \autoref{sec:e_architecture}). It runs in real time on a laptop CPU within the latency budget and reduces Gemini's word error rate from 53.3\,\% for the baseline to 39.6\,\%, while cutting the share of the interfering speaker's words that reach Gemini almost in half. On 123 held-out clips in which only the interfering speaker talks, it falls silent on 56 against 11 for WeSep, a larger offline system trained on different data, although WeSep recovers more of the target when both speakers talk (\autoref{sec:results_content_judge}).
    \item \textbf{A training loop with a proxy objective.} The live model cannot be differentiated through, thus the extractor is trained on a signal-level proxy objective adapted from CARTSE \cite{li2026cartse}.
  \end{enumerate}
  All results are for two-speaker mixtures and are optimised for Gemini.


The remainder of this report is structured as follows: \autoref{sec:litreview} reviews the background information and context for the target speaker extraction domain. \autoref{sec:methodology} describes the process of generating the data, the architecture of the extractor and the metrics. \autoref{sec:experiments} sets out how the model is trained, selected and evaluated. \autoref{sec:results} presents the results of the evaluation and \autoref{sec:conclusion} concludes the report with a summary and potential future work.